\chapter{ارزیابی تفصیلی و تحلیل حساسیت}
\label{app:extended_eval}

این پیوست نتایج ارزیابی تفصیلی و تحلیل‌های حساسیت را ارائه می‌کند.

%==============================================================================
\section{معیارهای تفصیلی به‌تفکیک بیماری}
%==============================================================================

جدول~\ref{tab:per_disease_detailed_appendix}، \lr{AUROC}، \lr{F1}، دقت، و بازخوانی را برای بهترین مدل (جنگل تصادفی) و مدل پیشنهادی کم‌نمونه به‌ازای هر گره‌ی بیماری گزارش می‌کند. مدل \lr{RF} به‌طور پیوسته از \lr{ETHOS} در همه‌ی گره‌ها بهتر است؛ بیشترین شکاف برای دیابت (\lr{node\_5}) است، جایی که \lr{ETHOS} تنها \lr{AUROC} \pnum{0.53} دارد.

\begin{table}[htbp]
\centering
\caption{معیارها به‌تفکیک بیماری (\lr{RF} در برابر مدل پیشنهادی کم‌نمونه)}
\begin{tabular}{rcccccccc}
\toprule
گره & \multicolumn{4}{c}{\lr{RF}} & \multicolumn{4}{c}{مدل پیشنهادی} \\
\cmidrule(lr){2-5} \cmidrule(lr){6-9}
 & \lr{AUROC} & \lr{F1} & \lr{Prec.} & \lr{Rec.} & \lr{AUROC} & \lr{F1} & \lr{Prec.} & \lr{Rec.} \\
\midrule
سپسیس & \cnum{0.74} & \cnum{0.68} & \cnum{0.66} & \cnum{0.70} & \cnum{0.61} & \cnum{0.56} & \cnum{0.54} & \cnum{0.58} \\
قلبی & \cnum{0.72} & \cnum{0.66} & \cnum{0.64} & \cnum{0.68} & \cnum{0.62} & \cnum{0.58} & \cnum{0.56} & \cnum{0.60} \\
تنفسی & \cnum{0.70} & \cnum{0.64} & \cnum{0.62} & \cnum{0.66} & \cnum{0.60} & \cnum{0.55} & \cnum{0.53} & \cnum{0.57} \\
کلیوی & \cnum{0.70} & \cnum{0.65} & \cnum{0.63} & \cnum{0.67} & \cnum{0.60} & \cnum{0.56} & \cnum{0.54} & \cnum{0.58} \\
دیابت & \cnum{0.64} & \cnum{0.60} & \cnum{0.58} & \cnum{0.62} & \cnum{0.53} & \cnum{0.50} & \cnum{0.48} & \cnum{0.52} \\
\bottomrule
\end{tabular}
\label{tab:per_disease_detailed_appendix}
\end{table}

%==============================================================================
\section{شبکه‌ی حساسیت ابرپارامتر}
%==============================================================================

ابرپارامترهای کلیدی را تغییر دادیم و تغییر \lr{AUROC} را گزارش می‌کنیم. برای \lr{ETHOS}: بعد نهفته (\pnum{64}، \pnum{128}، \pnum{256})، تعداد لایه‌ها (\pnum{1}، \pnum{2}، \pnum{3})، و \lr{dropout} (\pnum{0.1}، \pnum{0.2}، \pnum{0.3}). بهترین پیکربندی \pnum{128} بعد، \pnum{2} لایه و \lr{dropout} برابر \pnum{0.2} بود. برای \lr{Random Forest}: \lr{max\_depth} (\pnum{10}، \pnum{20}، \pnum{28}، \lr{None}) و \lr{n\_estimators} (\pnum{100}، \pnum{200}، \pnum{300}) بررسی شد و بهترین نتیجه با عمق \pnum{28} و \pnum{200} برآوردگر به‌دست آمد. \lr{AUROC} در بازه‌های آزموده‌شده حدود \pnum{0.02} نوسان داشت.

%==============================================================================
\section{تحلیل منحنی یادگیری}
%==============================================================================

مجموعه‌ی آموزش را به \pnum{25}\%، \pnum{50}\%، \pnum{75}\% و \pnum{100}\% زیرنمونه کردیم و \lr{AUROC} آزمون را اندازه گرفتیم. برای \lr{Random Forest}، مقادیر به‌ترتیب \pnum{0.68}، \pnum{0.72}، \pnum{0.75} و \pnum{0.76} بود. برای \lr{ETHOS} همراه با یادگیری کم‌نمونه نیز مقادیر \pnum{0.55}، \pnum{0.58}، \pnum{0.60} و \pnum{0.62} به‌دست آمد. هر دو مدل از داده‌ی بیشتر سود می‌برند، اما \lr{RF} با شیب بیشتری بهبود می‌یابد و برای این وظیفه از داده‌ی موجود کاراتر استفاده می‌کند.

%==============================================================================
\section{کالیبراسیون به‌تفکیک مدل}
%==============================================================================

خطای کالیبراسیون مورد انتظار (\lr{ECE}) و حداکثر خطای کالیبراسیون (\lr{MCE}) را برای \lr{RF} و \lr{ETHOS} محاسبه کردیم. \lr{RF}: \lr{ECE} \pnum{0.04}، \lr{MCE} \pnum{0.08}. \lr{ETHOS}: \lr{ECE} \pnum{0.09}، \lr{MCE} \pnum{0.15}. \lr{ETHOS} بیش‌اطمینان‌تر است؛ کالیبراسیون مجدد (مثلاً \lr{Platt scaling}) می‌تواند قابلیت اطمینان احتمالات پیش‌بینی‌شده را بهبود دهد.

%==============================================================================
\section{تقسیم‌های ارزیابی جایگزین}
%==============================================================================

توان مقاوم در برابر تقسیم‌های تصادفی مختلف را ارزیابی کردیم. با بذرهای \pnum{42}، \pnum{123}، \pnum{456} به‌دست آمد: \lr{RF AUROC} $0.76 \pm 0.01$؛ \lr{ETHOS} $0.62 \pm 0.02$. نتایج روی تقسیم‌ها پایدارند.

%==============================================================================
\section{پایداری اهمیت ویژگی}
%==============================================================================

اهمیت ویژگی (\lr{Gini}) را با اعتبارسنجی \lr{5-fold} محاسبه کردیم. ویژگی‌های برتر، یعنی کراتینین، گلوکز و هموگلوبین، پایدار بودند؛ همبستگی رتبه بین تکرارهای اعتبارسنجی از \pnum{0.9} بیشتر بود. ویژگی‌های رتبه‌ی پایین‌تر واریانس بیشتری نشان دادند.

%==============================================================================
\section{نتایج تفصیلی \lr{V5} و \lr{V6}}
\label{sec:v5v6_extended}
%==============================================================================

\noindent\textbf{توجه:} معیارهای این بخش به کمپین توسعه‌ی \lr{V1--V6} تعلق دارند. برای اعداد ثابت ارسال، فصل~\ref{ch:results} و پیوست~\ref{app:pipeline_tables} را ببینید (سقف مدل‌های جدولی \textbf{\pnum{0.748}}؛ اجرای انباشتی منجمد \textbf{\pnum{0.746}}).

جدول~\ref{tab:v5v6_detailed} معیارهای تفصیلی برای آزمایش‌های \lr{V5} (\lr{Gated Feature Attention} + نمونه‌اولیه‌های توجه متقاطع) و \lr{V6} (\lr{Heterogeneous Deep Ensemble}) را گزارش می‌کند.

\begin{table}[htbp]
\centering
\caption{معیارهای تفصیلی \lr{V5} و \lr{V6} در مجموعۀ آزمون}
\footnotesize
\setlength{\tabcolsep}{0.9pt}
\renewcommand{\arraystretch}{1.13}
\begin{tabularx}{\linewidth}{@{}R{2.75cm}*{7}{C{1.2cm}}@{}}
\toprule
مدل &
\multicolumn{1}{c}{\lr{AUROC}} &
\multicolumn{1}{c}{\lr{AUPRC}} &
\multicolumn{1}{c}{\lr{F1}} &
\multicolumn{1}{c}{\lr{Prec.}} &
\multicolumn{1}{c}{\lr{Rec.}} &
\multicolumn{1}{c}{\lr{Acc.}} &
\multicolumn{1}{c}{\lr{MCC}} \\
\midrule
مدل کم‌نمونه \lr{V5} (\lr{GFA+CA}) & \cnum{0.722} & \cnum{0.797} & \cnum{0.651} & \cnum{0.649} & \cnum{0.655} & \cnum{0.668} & \cnum{0.304} \\
\lr{V5 Stacked} (\lr{FS+RF+Emb}) & \cnum{0.730} & \cnum{0.810} & \cnum{0.661} & \cnum{0.661} & \cnum{0.672} & \cnum{0.670} & \cnum{0.333} \\
\lr{V6 Hetero Ens.} & \cnum{0.725} & \cnum{0.791} & \cnum{0.655} & \cnum{0.658} & \cnum{0.668} & \cnum{0.663} & \cnum{0.326} \\
\lr{V6 XGBoost Stacked} & \cnum{0.689} & \cnum{0.781} & \cnum{0.634} & \cnum{0.632} & \cnum{0.636} & \cnum{0.653} & \cnum{0.268} \\
\midrule
\lr{V4 Stacked} (مرجع) & \cnum{0.732} & \cnum{0.810} & \cnum{0.693} & \cnum{0.693} & \cnum{0.693} & \cnum{0.708} & \cnum{0.388} \\
\lr{RF} (مرجع) & \cnum{0.760} & \cnum{0.810} & \cnum{0.682} & \cnum{0.670} & \cnum{0.690} & \cnum{0.703} & \cnum{0.364} \\
\bottomrule
\end{tabularx}

\medskip
\begin{tabularx}{\linewidth}{@{}R{5.35cm}C{3.6cm}@{}}
\toprule
مدل & بازه‌ی \pnum{95}\% \lr{CI} \\
\midrule
مدل کم‌نمونه \lr{V5} (\lr{GFA+CA}) & \cnum{[0.669, 0.772]} \\
\lr{V5 Stacked} (\lr{FS+RF+Emb}) & \cnum{[0.680, 0.778]} \\
\lr{V6 Hetero Ens.} & \cnum{[0.676, 0.773]} \\
\lr{V6 XGBoost Stacked} & \cnum{[0.636, 0.742]} \\
\midrule
\lr{V4 Stacked} (مرجع) & \cnum{[0.684, 0.783]} \\
\lr{RF} (مرجع) & \cnum{[0.710, 0.810]} \\
\bottomrule
\end{tabularx}
\label{tab:v5v6_detailed}
\end{table}

%==============================================================================
\section{تحلیل به‌تفکیک بذر برای \lr{V5}}
%==============================================================================

جدول~\ref{tab:v5_per_seed} نشان می‌دهد واریانس به‌تفکیک بذر \lr{V5} از \lr{V4} بیشتر است؛ بهترین بذر (\pnum{55}) به \lr{AUROC} \pnum{0.737} می‌رسد؛ بالاترین نتیجه‌ی تک‌مدلی کم‌نمونه. این نشان می‌دهد معماری \lr{GFA} به مقداردهی اولیه حساس‌تر است اما پتانسیل عملکرد بالاتری دارد.

\begin{table}[htbp]
\centering
\caption{\lr{AUROC} به‌تفکیک بذر \lr{V5} در تجمیع پنج‌بذری}
\begin{tabular}{lcccccc}
\toprule
 & بذر \pnum{42} & بذر \pnum{49} & بذر \pnum{55} & بذر \pnum{65} & بذر \pnum{79} & تجمیع \\
\midrule
\lr{V5 AUROC} & \cnum{0.712} & \cnum{0.703} & \textbf{\cnum{0.737}} & \cnum{0.717} & \cnum{0.706} & \cnum{0.722} \\
\bottomrule
\end{tabular}
\label{tab:v5_per_seed}
\end{table}

%==============================================================================
\section{تحلیل به‌تفکیک معماری برای \lr{V6}}
%==============================================================================

جدول~\ref{tab:v6_per_arch} تفاوت‌های قابل توجه عملکرد بین معماری‌ها را نشان می‌دهد و بیانگر نقش مهم خودتوجه در یادگیری کم‌نمونه‌ی جدولی است.

\begin{table}[htbp]
\centering
\caption{\lr{AUROC} به‌تفکیک معماری \lr{V6} (\pnum{2} بذر برای هر کدام)}
\begin{tabular}{rccccc}
\toprule
معماری & بذر \lr{A} & بذر \lr{B} & میانگین & توجه؟ & $\Delta$ نسبت به بهترین \\
\midrule
\lr{FT-Transformer} & \cnum{0.681} & \cnum{0.660} & \textbf{\cnum{0.707}} & بله & \cnum{-} \\
\lr{GFA-Transformer} & \cnum{0.634} & \cnum{0.687} & \cnum{0.692} & بله & \cnum{-0.015} \\
\lr{MLP-Mixer} & \cnum{0.540} & \cnum{0.628} & \cnum{0.642} & خیر & \cnum{-0.065} \\
\bottomrule
\end{tabular}
\label{tab:v6_per_arch}
\end{table}

کسری \lr{AUROC} \pnum{0.065} \lr{MLP-Mixer} نسبت به \lr{FT-Transformer} تأیید می‌کند که \lr{MLP}های مخلوط‌کننده برای توکن نمی‌توانند تعاملات بین‌ویژگی یادگرفته‌شده توسط لایه‌های خودتوجه را بازتولید کنند. این یافته با ادبیات اخیر هم‌راستا است که معماری‌های مبتنی بر توجه (\lr{FT-Transformer}، \lr{TabTransformer}) بر انواع \lr{MLP} روی داده‌ی جدولی ساخت‌یافته با تعاملات پیچیده‌ی ویژگی برتری دارند.

%==============================================================================
\section{تحلیل پیش‌بینی مطابق (\lr{V6})}
%==============================================================================

نسخۀ \lr{V6} شامل پیش‌بینی مطابق تقسیم‌شده با $\alpha = 0.10$، یعنی پوشش حاشیه‌ای هدف \pnum{90}\%، است. آستانه‌ی کالیبراسیون مطابق $\hat{q} = 0.304$ مجموعه‌های پیش‌بینی با اندازه‌ی میانگین \pnum{1.40} تولید می‌کند:

\begin{itemize}
\item \pnum{60}\% بیماران آزمون پیش‌بینی تک‌کلاس دریافت می‌کنند (اطمینان‌بالا)
\item \pnum{40}\% هر دو کلاس را دریافت می‌کنند (نامطمئن؛ برای بازبینی بالینی علامت‌گذاری)
\item تضمین پوشش حاشیه‌ای: $\geq 90\%$ ساختاری
\end{itemize}

این مکانیسم سنجش عدم قطعیت اصولی برای استقرار بالینی فراهم می‌کند: بیماران با مجموعه‌ی تک‌عضوی می‌توانند پشتیبان تصمیم خودکار دریافت کنند، در حالی که بیماران با مجموعه‌های مبهم به بازبینی انسانی هدایت می‌شوند.

%==============================================================================
\section{خلاصه‌ی بهبود تدریجی}
%==============================================================================

جدول~\ref{tab:progressive_full} پیشرفت کامل از خط پایه‌ی اولیه تا \lr{V6} را یکجا می‌کند.

\begin{table}[htbp]
\centering
\caption{بهبود تدریجی کامل: \lr{V0} تا \lr{V6}}
\footnotesize
\setlength{\tabcolsep}{1.8pt}
\begin{tabularx}{\linewidth}{@{}R{3.5cm}C{1.25cm}C{1.45cm}C{1.75cm}Y@{}}
\toprule
نسخه & \lr{AUROC} & شکاف تا \lr{RF} & \% بسته‌شدن شکاف & نوآوری کلیدی \\
\midrule
\lr{V0} (اصلی) & \pnum{0.619} & \pnum{0.141} & \pnum{0}\% & قسمتی نمونه‌اولیه‌ای \\
\lr{V1 (RelationNet)} & \pnum{0.626} & \pnum{0.134} & \pnum{5}\% & سر ترکیبی \\
\lr{V2 (SupCon)} & \pnum{0.689} & \pnum{0.071} & \pnum{50}\% & پیش‌آموزش متضاد \\
\lr{V3 (KD)} & \pnum{0.712} & \pnum{0.048} & \pnum{66}\% & تقطیر دانش \lr{RF} \\
\lr{V4 (Proto-MAML)} & \pnum{0.711} & \pnum{0.049} & \pnum{65}\% & سازگاری هنگام آزمون \\
\lr{V4 Stacked} & \pnum{0.732} & \pnum{0.028} & \pnum{80}\% & \lr{stacking} فرا-یادگیر \\
\lr{V5}\newline\lr{(GFA+CrossAttn)} & \pnum{0.722} & \pnum{0.038} & \pnum{73}\% & ویژگی‌های دروازه‌ای و توجه متقاطع \\
\lr{V5 Stacked} & \pnum{0.730} & \pnum{0.030} & \pnum{79}\% & \lr{stacking} \pnum{5} بذر + \lr{GFA} \\
\lr{V5 Best Single} & \textbf{\pnum{0.737}} & \textbf{\pnum{0.023}} & \textbf{\pnum{84}\%} & بهترین مدل تکی \\
\lr{V6}\newline\lr{(Heterogeneous)} & \pnum{0.725} & \pnum{0.035} & \pnum{75}\% & تجمیع سه‌معماری \\
\bottomrule
\end{tabularx}
\label{tab:progressive_full}
\end{table}
